\documentclass[11pt,letterpaper]{article}
\usepackage[margin=1in]{geometry}
\usepackage{booktabs}
\usepackage{enumitem}
\usepackage{amssymb}
\usepackage[T1]{fontenc}
\usepackage{xcolor}
\usepackage{titlesec}
\usepackage{parskip}

\titleformat{\section}{\large\bfseries}{}{0em}{}[\titlerule]
\titleformat{\subsection}{\normalsize\bfseries}{}{0em}{}

\newcommand{\ok}{$\checkmark$}
\newcommand{\score}[1]{\textbf{#1}}

\title{\vspace{-2em}Grading Notes: \emph{Screen Time, Two Deletions}\\
\large AP Statistics $\cdot$ Unit 1 $\cdot$ Topics 1.1--1.7 $\cdot$ Bonus sheet}
\author{Second-read notes against the teacher key}
\date{September 23, 2026}

\begin{document}
\maketitle
\vspace{-2em}

\noindent\emph{Per the key, part (d) is hand-scored E/P/I. The scores below are a second read for you to confirm; borderline cases are flagged.}

\section{Summary}

\begin{center}
\begin{tabular}{@{}llccccl@{}}
\toprule
Student & Per. & (a) & (b) & (c) & (d) & Note \\
\midrule
Damian Cepeda        & E & \ok & \ok & \ok & \score{E} & P if strict on e3 \\
Yolandys Garcia Reyes& B & partial & \ok & \ok & \score{P+} & E if lenient on e2 \\
Zabdy Chevez         & B & \ok & \ok & \ok & \score{P} & \\
Rubiat Rahman        & B & \ok & \ok & mostly & \score{P} & \\
Elmer Lucero y Lucero& --- & partial & error & \ok & \score{P} & period blank \\
Mr.\ Co--- (unclear)  & B & partial & messy & \ok & \score{P} & \\
Tan Nguyen           & E & partial & \ok & \ok & \score{I} & closest to P \\
Cullen Freti         & B & partial & \ok & \ok & \score{I} & \\
Saly Ung             & B & partial & \ok & \ok & \score{I} & \\
Jesselly Ampudia     & E & \ok & \ok & blank & \score{I} & not attempted \\
Darla Gonzalez       & E & \ok & partial & partial & \score{I} & not attempted \\
\bottomrule
\end{tabular}
\end{center}

\section{Per-student notes}

\subsection{Damian Cepeda (E) --- \score{E}}
Cleanest sheet in the stack. (a)--(c) all correct; (c) is textbook. (d) hits e1 (10 of 38, over 240 min daily, ``does not fully answer whether screen time is a concern''), e2 (median, resistant), and e4 (deletion reasonable, report should say so). e3 is not stated inside (d) itself, only through ``resistant''; E is justified since (c) carries it, but P is defensible.

\subsection{Yolandys Garcia Reyes (B) --- \score{P+}}
Strongest e1 in the class and the only student to add the self-report caveat. e3 and e4 are both explicit. The gap: she never chooses a center to report---she reports the 26\% and mentions median resistance without committing to it. By the letter that is P; E is arguable. (c) has a slip: ``the mean changed by only 2 minutes'' should be \emph{median}. (a) is missing individuals and the screen-time value; First Take is blank.

\subsection{Zabdy Chevez (B) --- \score{P}}
Best (a) (``a randomly selected student'') and best explanation of \emph{why} the median resists (depends only on ordering). (d) used the frame verbatim: e2 and e3 present, no context for the 26\% and no disclosure statement. (b) wrote 24 freshmen but computed 36\%.

\subsection{Rubiat Rahman (B) --- \score{P}}
Same shape as Zabdy: solid (a)/(b); (d) is e2 + e3 only. (c) explains the mechanism but does not point to what in the histogram (the 720/900 bars) caused the move.

\subsection{Elmer Lucero y Lucero (no period) --- \score{P}}
(a) individuals given as ``student screen time'' (that is a variable). (b) wrote 18 seniors $\rightarrow$ 16.7\%; should be 15 $\rightarrow$ 20\%. (c) good. (d) has e2 and e3, no e1, and the disclosure part is muddled (``the principal should keep the data; removing the two values\ldots''). Check the period.

\subsection{Mr.\ Co--- (B, last name unreadable) --- \score{P}}
Good First Take (spotted missing what/units). (b) is heavily crossed out but the conclusion (``20\% $<$ 36\%, nowhere near as hyperbolic as 9 vs 3'') is right. (d) gives the 26\% context and picks the median, but says ``40 students, median 174'' (174 is from the 38) and ``with no outliers,'' which conceals the deletion rather than disclosing it---one of the LOOK FOR items.

\subsection{Tan Nguyen (E) --- \score{I}}
(a) missing individuals; (b)/(c) fine. (d): ``26\% does mean yes, screen time does affect students'' is exactly the floating-number-to-conclusion move. No center chosen, no resistance; deletion justified but not disclosed. I, but the closest I to P.

\subsection{Cullen Freti (B) --- \score{I}}
(b)/(c) fine. (a) values are odd (``Samsung,'' ``80000 minutes'') though the types are right. (d) misreads the 26\% as ``students whose screen time ranges from 0--600'' and trails off mid-sentence.

\subsection{Saly Ung (B) --- \score{I}}
Nearly identical to Cullen's (d): same 0--600 misreading, same trailing off; ``Nokia''/``1st Grade'' in (a). The two sheets resemble each other. (b) and (c) are good.

\subsection{Jesselly Ampudia (E) --- \score{I} (not attempted)}
(a) and (b) fully correct, then ``I don't get it'' for (c) and (d). She climbed two rungs of the ladder; a short conversation about resistance would likely unlock (d).

\subsection{Darla Gonzalez (E) --- \score{I} (not attempted)}
(a) correct. (b) blanks filled but no explanation sentence. (c) names the median and the two values but restates the prompt's numbers rather than explaining why the mean moved. (d) and First Take blank.

\section{Patterns worth noting}

\begin{enumerate}[leftmargin=*]
  \item \textbf{e1 and e4 are the most-missed elements}, while e2/e3 show up reliably. This matches the frame: ``I would report \rule{1cm}{0.4pt} because \rule{1cm}{0.4pt}; removing the two values changed \rule{1cm}{0.4pt} but not \rule{1cm}{0.4pt}'' steers students to e2 and e3 and never prompts for who/what/units or for saying the deletion happened. If the frame should scaffold all four elements, it needs two more slots.
  \item \textbf{Three Period B students read the 26\% as the share of students in the 0--600 range} rather than the share over 240. The ``10 of 38 are over 240'' sentence may deserve a moment in class.
\end{enumerate}

\end{document}
